\documentclass[17pt]{flash}
\usepackage{texmaths-tikz}
\def\figureA{%
    \begin{tikzpicture}[scale=\figScale, baseline=(current bounding box.center)]
        \Angle[angle=35,color=\currentcolor]
    \end{tikzpicture}
}
\def\figureB{%
    \begin{tikzpicture}[scale=\figScale, baseline=(current bounding box.center)]
        \Angle[angle=90,color=\currentcolor]
    \end{tikzpicture}
}
\def\figureC{%
    \begin{tikzpicture}[scale=\figScale, baseline=(current bounding box.center)]
        \Angle[angle=145,color=\currentcolor]
    \end{tikzpicture}
}
\def\figureD{%
    \begin{tikzpicture}[scale=\figScale, baseline=(current bounding box.center)]
        \Angle[angle=180,color=\currentcolor]
    \end{tikzpicture}
}
\startdatakeys
\cycle{3}
\connaissancesRequises{}
\competencesTravaillees{}
\enddatakeys
\begin{document}

\flashvisibletimertrue
\FlashSetItemSep{30pt}
\FlashSetVtop{0em}
\FlashSetTitle{Nature des angles}

\FlashShowConsignes{%
    \item Pour chaque angle, écrire sa nature.
}

\def\figScale{1}
\AtCorrection{\def\figScale{0.8}}
\flashlist{
    \flashitem{20}{\figureA}
    \flashitem{20}{\figureB}
    \flashitem{20}{\figureC}
    \flashitem{20}{\figureD}
}

\end{document}
